\documentclass[10pt,a4paper]{amsart}
\usepackage{amsmath,amssymb,amsthm,graphicx}
\usepackage{verbatim}
\newenvironment{reference}{\comment}{\endcomment}
%\newenvironment{reference}{}{}
\newenvironment{slogan}{\comment}{\endcomment}
\newenvironment{history}{\comment}{\endcomment}

% For commutative diagrams we use Xy-pic
\usepackage[all]{xy}

% We use 2cell for 2-commutative diagrams.
\xyoption{2cell}
% Theorem environments.
%
\theoremstyle{plain}
\newtheorem{theorem}[subsection]{Theorem}
\newtheorem{proposition}[subsection]{Proposition}
\newtheorem{lemma}[subsection]{Lemma}

\theoremstyle{definition}
\newtheorem{definition}[subsection]{Definition}
\newtheorem{example}[subsection]{Example}
\newtheorem{exercise}[subsection]{Exercise}
\newtheorem{situation}[subsection]{Situation}

\theoremstyle{remark}
\newtheorem{remark}[subsection]{Remark}
\newtheorem{remarks}[subsection]{Remarks}

\numberwithin{equation}{subsection}

% Macros
%
\def\lim{\mathop{\mathrm{lim}}\nolimits}
\def\colim{\mathop{\mathrm{colim}}\nolimits}
\def\Spec{\mathop{\mathrm{Spec}}}
\def\Hom{\mathop{\mathrm{Hom}}\nolimits}
\def\Ext{\mathop{\mathrm{Ext}}\nolimits}
\def\SheafHom{\mathop{\mathcal{H}\!\mathit{om}}\nolimits}
\def\SheafExt{\mathop{\mathcal{E}\!\mathit{xt}}\nolimits}
\def\Sch{\mathit{Sch}}
\def\Mor{\mathop{\mathrm{Mor}}\nolimits}
\def\Ob{\mathop{\mathrm{Ob}}\nolimits}
\def\Sh{\mathop{\mathit{Sh}}\nolimits}
\def\NL{\mathop{N\!L}\nolimits}
\def\CH{\mathop{\mathrm{CH}}\nolimits}
\def\proetale{{pro\text{-}\acute{e}tale}}
\def\etale{{\acute{e}tale}}
\def\QCoh{\mathit{QCoh}}
\def\Ker{\mathop{\mathrm{Ker}}}
\def\Im{\mathop{\mathrm{Im}}}
\def\Coker{\mathop{\mathrm{Coker}}}
\def\Coim{\mathop{\mathrm{Coim}}}

% Boxtimes
%
\DeclareMathSymbol{\boxtimes}{\mathbin}{AMSa}{"02}

%
% Macros for moduli stacks/spaces
%
\def\QCohstack{\mathcal{QC}\!\mathit{oh}}
\def\Cohstack{\mathcal{C}\!\mathit{oh}}
\def\Spacesstack{\mathcal{S}\!\mathit{paces}}
\def\Quotfunctor{\mathrm{Quot}}
\def\Hilbfunctor{\mathrm{Hilb}}
\def\Curvesstack{\mathcal{C}\!\mathit{urves}}
\def\Polarizedstack{\mathcal{P}\!\mathit{olarized}}
\def\Complexesstack{\mathcal{C}\!\mathit{omplexes}}
% \Pic is the operator that assigns to X its picard group, usage \Pic(X)
% \Picardstack_{X/B} denotes the Picard stack of X over B
% \Picardfunctor_{X/B} denotes the Picard functor of X over B
\def\Pic{\mathop{\mathrm{Pic}}\nolimits}
\def\Picardstack{\mathcal{P}\!\mathit{ic}}
\def\Picardfunctor{\mathrm{Pic}}
\def\Deformationcategory{\mathcal{D}\!\mathit{ef}}
\usepackage{cleveref}
\crefname{originallemma}{original lemma}{original lemmas}
\Crefname{originallemma}{Original lemma}{Original lemmas}
\makeatletter
\newlabel{lemma-cartesian-cocartesian}{{08N4}{source}{}{sourcepointer-1}{}}
\newlabel{lemma-cartesian-cocartesian@cref}{{[originallemma][0][]08N4}{[1][0][]source}{}{}{}}
\makeatother
\crefname{originallemma}{original lemma}{original lemmas}
\Crefname{originallemma}{Original lemma}{Original lemmas}
\makeatletter
\newlabel{lemma-cartesian-kernel}{{08N3}{source}{}{sourcepointer-2}{}}
\newlabel{lemma-cartesian-kernel@cref}{{[originallemma][0][]08N3}{[1][0][]source}{}{}{}}
\makeatother
\crefname{originallemma}{original lemma}{original lemmas}
\Crefname{originallemma}{Original lemma}{Original lemmas}
\makeatletter
\newlabel{lemma-check-exactness-fibre-product}{{08N5}{source}{}{sourcepointer-3}{}}
\newlabel{lemma-check-exactness-fibre-product@cref}{{[originallemma][0][]08N5}{[1][0][]source}{}{}{}}
\makeatother
\crefname{originallemma}{original lemma}{original lemmas}
\Crefname{originallemma}{Original lemma}{Original lemmas}
\makeatletter
\newlabel{lemma-exact-kernel-sequence}{{08N6}{source}{}{sourcepointer-4}{}}
\newlabel{lemma-exact-kernel-sequence@cref}{{[originallemma][0][]08N6}{[1][0][]source}{}{}{}}
\makeatother
\title{The snake lemma in an abelian category}
\author{The Stacks Project Authors}
\date{stacks-project:homology.tex:local-snapshot-20261012; GNU FDL 1.2 or later}
\begin{document}
\maketitle
\noindent\textit{Editorial note.} All delivered mathematical text is the author's own, copied from the pinned Stacks Project chapter file homology.tex; no statement or proof step was rewritten, reordered or removed.\par\medskip
\noindent\textit{Editorial note.} The proof is the author's complete argument. The closing sentence 'Then, the claim follows by Lemma lemma-exact-kernel-sequence and duality' names a specific same-chapter lemma, carried on this card as a navigation stub, and duality, rather than declining a step; all the computations it rests on are displayed inside the counted span. No sentence of the proof is a placeholder.\par\medskip
\noindent\textit{Editorial note.} Every labelled reference occurring inside the counted spans is carried in assembly.external\_references as a non-load-bearing navigation pointer (same-chapter ones at their definition site, cross-chapter ones at the invocation inside the counted span), so no load-bearing mathematics is replaced by a stub. No bibliography citation occurs inside the counted spans.\par\medskip
\noindent\textit{Editorial note.} The same-name check for this card scanned every staged Stacks chapter of this repository -- all sources/stacks-*/pinned-original-source.tex files, the card's own chapter excluded -- and the labels found are recorded in dedup.same\_name\_cross\_chapter on the card.\par\medskip
\noindent\textit{Editorial note.} The preamble is the chapter's own native preamble excerpt already used for this source: the comment-type reference, slogan and history environments (native preamble lines 8-13), the Xy-pic setup (native preamble lines 14-18) and the native theorem declarations with the mathematical macros (native preamble lines 154-217, the actual last line of the file).\par\medskip
\noindent\textit{Editorial note.} Same-name check: the label lemma-snake also occurs in the staged chapter stacks-algebra (sources/stacks-algebra/pinned-original-source.tex, line 521). The two are not the same claim and do not double count: the algebra version is the classical snake lemma for a diagram of abelian GROUPS and its proof ends with 'Proof of exactness is omitted.', while this card counts the author's abelian-CATEGORY version, which is more general and whose proof is written out in full inside the counted span. The more general statement is the one counted here, and the special case on the other side carries no proof of its own.\par\medskip
\noindent\textit{Editorial note.} Same-name check: the label lemma-snake-natural in the same chapter is a different, adjacent lemma (source line 1347) and is not part of this card's spans.\par\medskip
\noindent\textit{Allowed standard prerequisites (editorial).}\begin{itemize}
\item Cartesian and cocartesian squares in an abelian category. Same chapter, Lemma lemma-cartesian-cocartesian. Supplies the epimorphism and monomorphism facts about the canonical projections and coprojections that give uniqueness of the connecting map and the commutativity of its defining diagram.
\item Kernels from cartesian squares. Same chapter, Lemma lemma-cartesian-kernel. Produces the epimorphism from x onto the kernel of the projection pi used in the construction of the connecting map.
\item Testing exactness through a fibre product. Same chapter, Lemma lemma-check-exactness-fibre-product. Converts the exactness required at ker(gamma) into the objects s and t and the morphisms n, zeta and vartheta that carry the rest of the argument.
\item Exactness of the kernel sequence. Same chapter, Lemma lemma-exact-kernel-sequence. Supplies the last step from the exactness of ker(beta) -> ker(gamma) -> coker(alpha) to the exactness of the full six-term sequence of the statement.
\end{itemize}\medskip
\ifdefined\hypertarget\hypertarget{sourcepointer-1}{}\fi
\noindent\textit{Source reference.} Original lemma 08N4: sources/stacks-homology/pinned-original-source.tex, source lines 1104--1104. Navigation pointer invoked at proof line 1174; the counted span contains the pointer only and the cited result is not reproduced inside it, so the pointer is registered as a non-load-bearing navigation stub in the established form for Stacks material. No bibliography citation is involved.\par\medskip
\ifdefined\hypertarget\hypertarget{sourcepointer-2}{}\fi
\noindent\textit{Source reference.} Original lemma 08N3: sources/stacks-homology/pinned-original-source.tex, source lines 1069--1069. Navigation pointer invoked at proof line 1292; the counted span contains the pointer only and the cited result is not reproduced inside it, so the pointer is registered as a non-load-bearing navigation stub in the established form for Stacks material. No bibliography citation is involved.\par\medskip
\ifdefined\hypertarget\hypertarget{sourcepointer-3}{}\fi
\noindent\textit{Source reference.} Original lemma 08N5: sources/stacks-homology/pinned-original-source.tex, source lines 1152--1152. Navigation pointer invoked at proof line 1216; the counted span contains the pointer only and the cited result is not reproduced inside it, so the pointer is registered as a non-load-bearing navigation stub in the established form for Stacks material. No bibliography citation is involved.\par\medskip
\ifdefined\hypertarget\hypertarget{sourcepointer-4}{}\fi
\noindent\textit{Source reference.} Original lemma 08N6: sources/stacks-homology/pinned-original-source.tex, source lines 1186--1186. Navigation pointer invoked at proof line 1342; the counted span contains the pointer only and the cited result is not reproduced inside it, so the pointer is registered as a non-load-bearing navigation stub in the established form for Stacks material. No bibliography citation is involved.\par\medskip
\begin{lemma}
\label{lemma-snake}
Let $\mathcal{A}$ be an abelian category. Let
$$
\xymatrix{
& x \ar[r]^f \ar[d]^\alpha &
y \ar[r]^g \ar[d]^\beta &
z \ar[r] \ar[d]^\gamma &
0 \\
0 \ar[r] & u \ar[r]^k & v \ar[r]^l & w
}
$$
be a commutative diagram with exact rows.
\begin{enumerate}
\item There exists a unique morphism
$\delta : \Ker(\gamma) \to \Coker(\alpha)$
such that the diagram
$$
\xymatrix{
y \ar[d]_\beta &
y \times_z \Ker(\gamma) \ar[l]_{\pi'} \ar[r]^{\pi} &
\Ker(\gamma) \ar[d]^\delta \\
v \ar[r]^{\iota'} & \Coker(\alpha) \amalg_u v &
\Coker(\alpha) \ar[l]_\iota
}
$$
commutes, where $\pi$ and $\pi'$ are the canonical projections
and $\iota$ and $\iota'$ are the canonical coprojections.
\item The induced sequence
$$
\Ker(\alpha) \xrightarrow{f'} \Ker(\beta) \xrightarrow{g'}
\Ker(\gamma) \xrightarrow{\delta} \Coker(\alpha) \xrightarrow{k'}
\Coker(\beta) \xrightarrow{l'} \Coker(\gamma)
$$
is exact. If $f$ is injective then so is $f'$, and if $l$ is
surjective then so is $l'$.
\end{enumerate}
\end{lemma}

\begin{proof}
As $\pi$ is an epimorphism and $\iota$ is a monomorphism by
Lemma \ref{lemma-cartesian-cocartesian}, uniqueness of $\delta$ is clear.
Let $p = y \times_z \Ker(\gamma)$ and $q = \Coker(\alpha) \amalg_u v$.
Let $h : \Ker(\beta) \to y$, $i : \Ker(\gamma) \to z$ and
$j : \Ker(\pi) \to p$ be the canonical injections.
Let $\pi'' : u \to \Coker(\alpha)$ be the canonical projection.
Keeping in mind Lemma \ref{lemma-cartesian-cocartesian} we get a commutative
diagram with exact rows
$$
\xymatrix{
0 \ar[r] &
\Ker(\pi) \ar[r]^j &
p \ar[r]^{\pi} \ar[d]_{\pi'} &
\Ker(\gamma) \ar[d]_i \ar[r] & 0 \\
& x \ar[r]^f \ar[d]_\alpha & y \ar[r]^g \ar[d]_\beta &
z \ar[d]_\gamma \ar[r] & 0 \\
0 \ar[r] & u \ar[r]^k \ar[d]_{\pi''} &
v \ar[r]^l \ar[d]_{\iota'} & w & \\
0 \ar[r] & \Coker(\alpha) \ar[r]^\iota & q & &
}
$$
As $l \circ \beta \circ \pi' = \gamma \circ i \circ \pi = 0$ and as the third
row of the diagram above is exact, there is an $a : p \to u$
with $k \circ a = \beta \circ \pi'$. As the upper right quadrangle of the
diagram above is cartesian, Lemma \ref{lemma-cartesian-kernel} yields an
epimorphism $b : x \to \Ker(\pi)$ with $\pi' \circ j \circ b = f$.
It follows
$k \circ a \circ j \circ b = \beta \circ \pi' \circ j \circ b =
\beta \circ f = k \circ \alpha$.
As $k$ is a monomorphism this implies $a \circ j \circ b = \alpha$. It follows
$\pi'' \circ a \circ j \circ b = \pi'' \circ \alpha = 0$. As $b$ is an
epimorphism this
implies $\pi'' \circ a \circ j = 0$. Therefore, as the top row of the diagram
above is exact, there exists
$\delta : \Ker(\gamma) \to \Coker(\alpha)$ with
$\delta \circ \pi = \pi'' \circ a$. It follows
$\iota \circ \delta \circ \pi = \iota \circ \pi'' \circ a =
\iota' \circ k \circ a = \iota' \circ \beta \circ \pi'$
as desired.

\medskip\noindent
As the upper right quadrangle in the diagram above is cartesian there
is a $c : \Ker(\beta) \to p$ with $\pi' \circ c = h$ and $\pi \circ c = g'$.
It follows
$\iota \circ \delta \circ g' = \iota \circ \delta \circ \pi \circ c =
\iota' \circ \beta \circ \pi' \circ c = \iota' \circ \beta \circ h = 0$.
As $\iota$ is a monomorphism this implies $\delta \circ g' = 0$.

\medskip\noindent
Next, let $d : r \to \Ker(\gamma)$ with $\delta \circ d = 0$. Applying
Lemma \ref{lemma-check-exactness-fibre-product} to the exact sequence
$p \xrightarrow{\pi} \Ker(\gamma) \to 0$ and $d$ yields an object $s$,
an epimorphism $m : s \to r$ and a morphism $n : s \to p$ with
$\pi \circ n = d \circ m$. As
$\pi'' \circ a \circ n = \delta \circ d \circ m = 0$,
applying Lemma \ref{lemma-check-exactness-fibre-product} to the exact sequence
$x \xrightarrow{\alpha} u \xrightarrow{p} \Coker(\alpha)$ and
$a \circ n$ yields an object $t$, an epimorphism $\varepsilon : t \to s$ and
a morphism $\zeta : t \to x$ with
$a \circ n \circ \varepsilon = \alpha \circ \zeta$.
It holds
$\beta \circ \pi' \circ n \circ \varepsilon =
k \circ \alpha \circ \zeta = \beta \circ f \circ \zeta$.
Let $\eta = \pi' \circ n \circ \varepsilon - f \circ \zeta : t \to y$. Then,
$\beta \circ \eta = 0$. It follows that there is a
$\vartheta : t \to \Ker(\beta)$ with $\eta = h \circ \vartheta$. It holds
$i \circ g' \circ \vartheta = g \circ h \circ \vartheta =
g \circ \pi' \circ n \circ \varepsilon - g \circ f \circ \zeta =
i \circ \pi \circ n \circ \varepsilon = i \circ d \circ m \circ \varepsilon$.
As $i$ is a monomorphism we get
$g' \circ \vartheta = d \circ m \circ \varepsilon$.
Thus, as $m \circ \varepsilon$ is an epimorphism,
Lemma \ref{lemma-check-exactness-fibre-product} implies that
$\Ker(\beta) \xrightarrow{g'} \Ker(\gamma) \xrightarrow{\delta} \Coker(\alpha)$
is exact. Then, the claim follows by Lemma \ref{lemma-exact-kernel-sequence}
and duality.
\end{proof}

\end{document}
